%Chargement des paquets
\usepackage{amsmath}
\usepackage{amsfonts}
\usepackage{amssymb}
\usepackage{enumerate}
\usepackage{mathtools}
\usepackage{bbm}
\usepackage{xparse, etoolbox}
\usepackage{enumerate}
\usepackage{enumitem}
\usepackage{mathabx}
\usepackage{babel}[french]

%environment
\usepackage{amsthm}
\usepackage{keytheorems}
\usepackage{hyperref} 

%Interface théorème
\newkeytheorem{lem}[
	name = Lemme
]

\newkeytheorem{prop}[
	name = Proposition
]

\newkeytheorem{thm}[
	name = Théorème
]

\newkeytheorem{coro}[
	name = Corollaire
]

\renewcommand*{\proofname}{Démonstration}

\theoremstyle{definition}
\newtheorem{defn}{Définition}

\theoremstyle{definition}
\newtheorem*{nota}{Notation}

\theoremstyle{definition}
\newtheorem*{limi}{Liminaire}

\theoremstyle{remark}
\newtheorem{rmq}{Remarque}

\theoremstyle{plain}
\newtheorem*{fait}{Fait}


%Ensembles de nombres
\newcommand{\N} {\mathbb{N}}
\newcommand{\Ne}{\N^\ast}
\newcommand{\Z} {\mathbb{Z}}
\newcommand{\Q} {\mathbb{Q}}
\newcommand{\R} {\mathbb{R}}
\newcommand{\Rb}{\overline{\mathbb{R}}}
\newcommand{\Rp}{\R_+}
\newcommand{\Rm}{\R_-}
\newcommand{\K} {\mathbb{K}}
\newcommand{\Ket} {\mathbb{K^{\ast}}}
\newcommand{\Cx}{\mathbb{C}}

%Ensembles, tribus
\newcommand{\A}{\mathbb{A}}
\renewcommand{\P}{\mathcal{P}}
\newcommand{\T}{\mathcal{T}}
\newcommand{\F}{\mathcal{F}}
\newcommand{\C} {\mathcal{C}}
\newcommand{\ouv}{\mathcal{O}}
\newcommand{\B}{\mathcal{B}}
\newcommand{\M}{\mathcal{M}}
\newcommand{\etag}{\mathcal{E}}
\newcommand{\etagOp}{\etag^{+}(\Omega)}
\newcommand{\etagO}{\etag(\Omega)}
%%Lp avant quotientage
\newcommand{\Lic}{\mathcal{L}^1}
\newcommand{\Lpc}{\mathcal{L}^p}
\newcommand{\Linfc}{\mathcal{L}^{\infty}}
\newcommand{\Lifull}{\Lic(\Omega, \B, m)}
\newcommand{\Lpfull}{\Lpc(\Omega, \B, m)}
\newcommand{\Linffull}{\Linfc(\Omega, \B, m)}
%%Lp après quotientage
\newcommand{\Li}{L^1}
\newcommand{\Ld}{L^2}
\newcommand{\Lp}{L^p}
\newcommand{\Linf}{L^{\infty}}
\newcommand{\Listd}{\Li(\Omega, m)} 
\newcommand{\Ldstd}{\Ld(\Omega, m)} 
\newcommand{\Lpstd}{\Lp(\Omega, m)} 

%Quelques opérateurs
\DeclareMathOperator{\codim}{codim}
\DeclareMathOperator{\rg}{rg}
\DeclareMathOperator{\tr}{tr}
\DeclareMathOperator{\vect}{Vect}
\DeclareMathOperator{\ima}{Im}
\DeclareMathOperator{\id}{Id}
\DeclareMathOperator{\sign}{sign}
\DeclareMathOperator{\ext}{ext}
\newcommand{\ind}{\mathbbm{1}}
\newcommand*{\spr}[2]{\left(\,#1\,\middle|\,#2\,\right)}
\newcommand{\interior}[1]{%
  {\kern0pt#1}^{\mathrm{o}}%
}
\DeclareMathOperator{\card}{card}
\DeclareMathOperator{\ppcm}{ppcm}

%Commandes de pure flemme
\newcommand{\veps}{\varepsilon}
\newcommand{\intab}{\int_{a}^{b}}
\newcommand{\intR}{\int_{-\infty}^{\infty}}
\newcommand{\intinf}{\int_{- \infty}^{+ \infty}}
\newcommand{\equi}{\Leftrightarrow}
\newcommand{\Kev}{$\K$-espace vectoriel }
\newcommand{\Kevs}{$\K$-espaces vectoriels }
\newcommand{\Rev}{$\R$-espace vectoriel }
\newcommand{\Revs}{$\R$-espaces vectoriels }
\newcommand{\Cev}{$\Cx$-espace vectoriel }
\newcommand{\Cevs}{$\Cx$-espaces vectoriels }
\newcommand{\evn}{(E, \norm{.})}
\newcommand{\interf}[2]{[#1,#2]}
\newcommand{\limb}{\lim\limits_}
\newcommand{\lebn}{\lambda_n}
\newcommand{\pp}{\text{~presque partout}}
\newcommand{\derivp}[2]{\frac{\partial #1}{\partial #2}}
\newcommand{\famiu}{(u_i)_{i \in I}}
\newcommand{\sev}{\text{~s.e.v.}}
\newcommand{\Mnp}{M_{n,p}(\K)}
\newcommand{\Mn}{M_{n}(\K)}
\newcommand{\tq}{\text{~t.q.~}}
\newcommand{\mq}{~montrer~que~}
\newcommand{\et}{\quad \text{et} \quad}
\newcommand{\ou}{\text{~ou~}}
\newcommand{\Kx}{\K[X]}


%Commandes de gars chelou
\renewcommand{\subset}{\subseteq}

%Commande restriction, ty duckduckgo
\def\restr#1#2{\mathchoice
              {\setbox1\hbox{${\displaystyle #1}_{\scriptstyle #2}$}
              \restraux{#1}{#2}}
              {\setbox1\hbox{${\textstyle #1}_{\scriptstyle #2}$}
              \restraux{#1}{#2}}
              {\setbox1\hbox{${\scriptstyle #1}_{\scriptscriptstyle #2}$}
              \restraux{#1}{#2}}
              {\setbox1\hbox{${\scriptscriptstyle #1}_{\scriptscriptstyle #2}$}
              \restraux{#1}{#2}}}
\def\restraux#1#2{{#1\,\smash{\vrule height .8\ht1 depth .85\dp1}}_{\,#2}} 

%Quelques normes
\DeclarePairedDelimiter\abs{\lvert}{\rvert}
\DeclarePairedDelimiter\labs{\bigl\lvert}{\bigr\rvert}
\DeclarePairedDelimiter\norm{\lVert}{\rVert}
\DeclarePairedDelimiterXPP\normi[1]{}\lVert\rVert{_1}{\ifblank{#1}{\:\cdot\:}{#1}}
\DeclarePairedDelimiterXPP\normd[1]{}\lVert\rVert{_2}{\ifblank{#1}{\:\cdot\:}{#1}}
\DeclarePairedDelimiterXPP\normp[1]{}\lVert\rVert{_p}{\ifblank{#1}{\:\cdot\:}{#1}}
\DeclarePairedDelimiterXPP\normq[1]{}\lVert\rVert{_q}{\ifblank{#1}{\:\cdot\:}{#1}}
\DeclarePairedDelimiterXPP\norminf[1]{}\lVert\rVert{_\infty}{\ifblank{#1}{\:\cdot\:}{#1}}

%Bracket en angle
\DeclarePairedDelimiter\angl{\langle}{\rangle}

%Set, parenthèses
\newcommand{\set}[1]{\{ #1 \}}
\newcommand{\bigset}[1]{\bigl\{ #1 \bigr\}}
\newcommand{\Bigset}[1]{\Bigl\{ #1 \Bigr\}}
\newcommand{\bigpar}[1]{\bigl( #1 \bigr)}
\newcommand{\Bigpar}[1]{\Bigl( #1 \Bigr)}

%Opitmisation des opérateurs de comparaison
\DeclarePairedDelimiter\tio{o (}{)}
\DeclarePairedDelimiter\gro{O (}{)}


\pagestyle{fancy}

\fancyhead[C]{Le théorème du point fixe de Bouwer}
\fancyhead[R]{}

\begin{document}

\underline{Source :} Rouvière - Calcul différentiel - page 165. \newline

\underline{Leçon :} 204 : Connexité. Exemples d’applications. \newline

\begin{nota}
On suppose $\R^n, n \geq 1$ muni du produit scalaire canonique, dénoté par $\angl{\cdot, \cdot}$, 
et de la norme induite $\norm{\cdot}$.
On note respectivement $B$ et $S$ la boule unité fermée et la sphère unité de $\R^n$ (pour cette norme).
\end{nota}

\begin{lem}
\label{lem:lacets}
Il n'existe pas de retraction forte par déformation de la boule unité fermée de $\R^n$ sur la sphère unité, 
c'est-à-dire d'application $G :  B \to S$ continue et dont la restriction à $S$ est l'identité. 
\end{lem}

\begin{proof}
Supponsons qu'une telle application $G$ existe, alors on définit la fonction $F$ par :
\[ \forall t,x \in [0,1], \quad F(t,x) = G \left ( xe^{i2\pi t} \right) , \]
Pour $x=1$ on obtient un lacet d'extrémité $F(0)$ et qui fait le tour du disque. 
Pour $x=0$, on obtient un lacet constant de valeur $F(0)$.
Or ces deux lacets sont homotopes, ce qui est absurde (vous pouvez me faire confiance, j'ai essayé à la maison).
\end{proof}

\begin{thm}[du point fixe de Brouwer]
Toute application continue de $B$ dans $B$ admet au moins un point fixe.
\end{thm}

\begin{proof}
Soit $f$ une telle fonction.

\begin{enumerate}[label=(\roman*)]

\item
On commence par démontrer le résultat pour $n=1$, c'est-à-dire pour une fonction $f$ continue de $[-1, 1]$ dans lui-même.
Supposons que $f$ n'admette pas de point fixe sur cet intervalle, par continuité on a :
\[ \forall x \in [-1, 1], \quad f(x) > x \ou f(x) < x , \]
et aucune de ces conditions n'est possible aux bornes de l'intervalle. Le résultat est donc démontré pour $n=1$.

\item
On fixe maintenant $n \geq 2$ et on suppose à nouveau le résultat faux. 
On va montrer qu'il existe alors une rétraction forte par déformation de la boule unité sur la sphère, que l'on notera $G$.
Soit donc $f : B \to B$ sans point fixe. 
Pour tout $x \in B , f(x) \neq x$, on peut donc définir une unique demi-droite $(D_x)$ d'origine $f(x)$ et passant par $x$.
En notant $G(x)$ l'intersection de cette demi-droite avec $S$ (origine exclue), on a défini une application de $B$ dans $S$. 
Il reste à montrer que cette application est continue et que sa restriction à $S$ est l'identité. \\
Le second point est clair, $x$ étant pris sur $S$.
On peut paramétrer la droite $(D_x)$ par $(D_x) : \lambda x +(1-\lambda) f(x)$ avec $\lambda >0$. 
Pour trouver l'intersection de la demi-droite avec $S$ on doit donc trouver $\lambda > 0$ tel que :
\[ \norm{\lambda x + (1-\lambda)f(x)}^2 = 1 , \]
ce qui équivaut à résoudre :
\[ P_x(\lambda) = \lambda^2 \norm{x-f(x)}^2 + 2 \lambda \angl{ x - f(x), f(x)} + \norm{f(x)}^2 - 1 = 0 . \]
On doit donc résoudre une équation du second degré (en $\lambda$) de discriminant 
\begin{align*}
\Delta(x) & = (2\angl{ x - f(x), f(x)})^2 - 4  \norm{x-f(x)}^2 (\norm{f(x)}^2 - 1) \\
& = 4 \Bigpar{ \angl{ x - f(x), f(x)}^2 - \norm{x-f(x)}^2 \norm{f(x)}^2 } + 4 \norm{x-f(x)}^2 \\
& = A + B ,
\end{align*}
avec $A \geq 0$ (c'est l'inégalité de Cauchy-Schwarz) et $B > 0$ ($f$ n'admet pas de point fixe). 
Donc $P_x$ admet deux racines réelles distinctes. 
Comme de plus $P_\lambda(0) \leq 0$ (et $\norm{x-f(x)}^2 > 0$),  $P_x$ admet nécessairement une racine sur $\Rp$. 
Cette racine a pour expression explicite :
\[ \lambda(x) = \dfrac{ - \angl{ x - f(x), f(x)} + \sqrt{\Delta(x)} }{ \norm{x-f(x)}^2 } , \]
qui est une fonction continue en $x$. Or $G(x) = \lambda(x)$, on a donc démontré la continuité de $G$. \\
Ainsi, on arrive à une contradiction selon le lemme $\ref{lem:lacets}$ ce qui nous permet de conclure.
\end{enumerate}

\end{proof}

\end{document}